\section*{1. Centered data and covariance}
Let $D\in\R^{n\times p}$ contain $n\geqslant2$ observations as rows and
$p$ variables as columns. Put
\[
\bar d=\frac1nD^\top\mathbf1\in\R^p,\qquad
C=I_n-\frac1n\mathbf1\mathbf1^\top,\qquad
Z=CD=D-\mathbf1\bar d^\top.
\]
Thus row $i$ of $Z$ is the deviation of observation $i$ from the mean
row. Since $\mathbf1^\top\mathbf1=n$,
$C^\top=C$, $C^2=C$, $C\mathbf1=0$ and $\mathbf1^\top Z=0$.
The sample covariance matrix is
\[
S=\frac1{n-1}Z^\top Z\in\R^{p\times p}.
\]
Its entry $s_{jk}$ is the sample covariance between columns $j,k$.
It is symmetric positive semidefinite because
$a^\top Sa=\norm{Za}^2/(n-1)\geqslant0$.
Moreover $\ker S=\ker Z$, so
\[
\operatorname{rank}S=\operatorname{rank}Z\leqslant\min(p,n-1),\qquad
\tr S=\frac{\norm Z_F^2}{n-1}=\sum_{j=1}^p s_{jj}.
\]
The rank bound uses $\operatorname{col}(Z)\subseteq\mathbf1^\perp$.
The total variance is the sum of the individual variable variances.

For $a\in\R^p$, the column $Za$ contains the centered values of the
linear combination with coefficients $a$. Its mean is zero and its
sample variance is $a^\top Sa$. A norm constraint on $a$ is essential:
replacing $a$ by $ta$ multiplies the variance by $t^2$.

\begin{example}[A complete centered covariance calculation]\label{eg:12:data}
Take
\[
D=\begin{pmatrix}12&21\\11&22\\9&18\\8&19\end{pmatrix}.
\]
The mean is $(10,20)^\top$, giving
\[
Z=\begin{pmatrix}2&1\\1&2\\-1&-2\\-2&-1\end{pmatrix},\qquad
Z^\top Z=\begin{pmatrix}10&8\\8&10\end{pmatrix},\qquad
S=\frac13\begin{pmatrix}10&8\\8&10\end{pmatrix}.
\]
For $a=e_1$, the score variance is $10/3$; for
$a=(1,1)^\top/\sqrt2$, it is $(10+8+8+10)/6=6$.
The unit-length direction changes which linear combination is being
measured while fixing the coefficient scale.
\end{example}

\section*{2. Principal directions and scores}
Write the spectral decomposition
\[
S=V\Lambda V^\top,\qquad V=(v_1\ \cdots\ v_p),\quad
\Lambda=\operatorname{diag}(\lambda_1,\ldots,\lambda_p),\quad
\lambda_1\geqslant\cdots\geqslant\lambda_p\geqslant0.
\]
The \textbf{principal directions}, or loading vectors in these notes, are
the unit columns $v_j$. The \textbf{score matrix} is
$T=ZV\in\R^{n\times p}$, with score column $t_j=Zv_j$.
The meaning of ``loading'' is fixed here as the unit eigenvector, not
an eigenvector multiplied by a square root of its eigenvalue.

\begin{theorem}[Successive maximum variance]\label{thm:12:successive}
Among unit $a$, the maximum of $a^\top Sa$ is $\lambda_1$.
After choosing orthonormal eigenvectors $v_1,\ldots,v_{k-1}$ in decreasing
eigenvalue order, the maximum under the additional restrictions
$a^\top v_i=0$ for $i<k$ is $\lambda_k$, attained at $v_k$.
The score columns have sample covariance
$T^\top T/(n-1)=\Lambda$.
\end{theorem}
\begin{proof}
Expand $a=\sum_i c_iv_i$. The unit constraint becomes $\sum_i c_i^2=1$,
and $a^\top Sa=\sum_i\lambda_i c_i^2$.
It is at most $\lambda_1$, with equality precisely when $c_i=0$ for
every $\lambda_i<\lambda_1$.
The extra orthogonality restrictions set $c_1=\cdots=c_{k-1}=0$.
The remaining weighted average is at most $\lambda_k$, attained by
putting $c_k=1$ and all other coefficients zero. Its equality vectors
are the unit vectors in the still-available part of the $\lambda_k$
eigenspace. Finally,
$T^\top T/(n-1)=V^\top SV=\Lambda$.
\end{proof}

Zero covariance between score columns follows from this matrix identity.
Statistical independence is a separate assertion and need not hold.
When $\tr S>0$, the variance proportion of component $j$ is
$\lambda_j/\tr S$, and the first $k$ components explain
$\sum_{j=1}^k\lambda_j/\tr S$ of total variance.
If $S=0$, all observations equal their mean and these proportions have
zero denominator; no direction has positive variance.

\begin{example}[Scores, reconstruction, and a residual]\label{eg:12:scores}
For Example~\ref{eg:12:data},
\[
\lambda_1=6,\quad \lambda_2=\tfrac23,\qquad
v_1=\tfrac1{\sqrt2}\begin{pmatrix}1\\1\end{pmatrix},\quad
v_2=\tfrac1{\sqrt2}\begin{pmatrix}1\\-1\end{pmatrix}.
\]
The scores are
\[
T=\frac1{\sqrt2}\begin{pmatrix}3&1\\3&-1\\-3&1\\-3&-1\end{pmatrix}.
\]
Their squared lengths are $18$ and $2$, and their dot product is zero.
Dividing by $n-1=3$ recovers variances $6$ and $2/3$.
The first direction explains $6/(6+2/3)=9/10$ of total variance.
Keeping only its score reconstructs
\[
Z_1=t_1v_1^\top=\frac12\begin{pmatrix}3&3\\3&3\\-3&-3\\-3&-3\end{pmatrix},\qquad
D_1=\mathbf1\bar d^\top+Z_1.
\]
For example the first reconstructed observation is $(11.5,21.5)$.
The residual is
\[
Z-Z_1=\frac12\begin{pmatrix}1&-1\\-1&1\\1&-1\\-1&1\end{pmatrix},
\]
whose squared Frobenius norm is $8/4=2$.
The lost variance is $2/(n-1)=2/3$, exactly the discarded eigenvalue.
\end{example}

\section*{3. Variance captured by a subspace}
Let $Q\in\R^{p\times k}$ have orthonormal columns.
Projection of each centered row onto $\operatorname{col}(Q)$ gives
$ZQQ^\top$. Its score coordinates in this basis are $ZQ$, and
\[
\frac{\norm{ZQ}_F^2}{n-1}=\tr(Q^\top SQ).
\]
Replacing $Q$ by $QR$ for an orthogonal $k\times k$ matrix $R$ rotates
the score coordinates but leaves $QQ^\top$ and the reconstruction unchanged.

\begin{theorem}[Maximum captured variance]\label{thm:12:trace}
For any symmetric $p\times p$ matrix $A$ with eigenvalues in decreasing
order and any $1\leqslant k\leqslant p$,
\[
\max_{Q^\top Q=I_k}\tr(Q^\top AQ)=\sum_{i=1}^k\lambda_i(A),\qquad
\min_{Q^\top Q=I_k}\tr(Q^\top AQ)=\sum_{i=p-k+1}^p\lambda_i(A).
\]
For the maximum, an optimal subspace contains every eigenspace with
eigenvalue strictly greater than $\lambda_k$ and is contained in the
sum of the eigenspaces with eigenvalue at least $\lambda_k$.
\end{theorem}
\begin{proof}
Let $A=V\Lambda V^\top$ and $P=QQ^\top$. Put $w_i=v_i^\top Pv_i$.
Since $P$ is an orthogonal projector, $w_i=\norm{Pv_i}^2\in[0,1]$ and
$\sum_i w_i=\tr P=k$. Cyclic trace gives
$\tr(Q^\top AQ)=\sum_i\lambda_iw_i$.
Set $\alpha=\lambda_k$. Then
\[
\begin{split}
\sum_{i=1}^k\lambda_i-\sum_{i=1}^p\lambda_iw_i
={}&\sum_{i\leqslant k}(\lambda_i-\alpha)(1-w_i)
+\sum_{i>k}(\alpha-\lambda_i)w_i\geqslant0.
\end{split}
\]
Here the terms in $\alpha$ cancel because $\sum_iw_i=k$.
The upper bound is attained by $Q=(v_1\ \cdots\ v_k)$.
Equality forces $w_i=1$ when $\lambda_i>\alpha$ and $w_i=0$ when
$\lambda_i<\alpha$. The identities $w_i=1\iff Pv_i=v_i$ and
$w_i=0\iff Pv_i=0$ give the stated subspace characterization, including
ties. Conversely those conditions give equality.
Apply the upper-bound result to $-A$ to obtain the minimum, attained
by the bottom $k$ eigenvectors.
\end{proof}

If $k<p$ and $\lambda_k>\lambda_{k+1}$, the optimal $k$-dimensional
subspace is unique, although its orthonormal basis is not.
If the cutoff lies inside a repeated eigenspace, several optimal
subspaces may exist. Changing the sign of a loading also changes the
sign of its score, leaving their product $t_jv_j^\top$ unchanged.

\section*{4. Best low-rank reconstruction}
The Frobenius norm satisfies
$\norm{A}_F^2=\sum_{i,j}a_{ij}^2=\tr(A^\top A)$ and is unchanged by
orthogonal multiplication on either side. If $P$ is an orthogonal
projector, rowwise Pythagoras gives
\[
\norm Z_F^2=\norm{ZP}_F^2+\norm{Z(I-P)}_F^2.
\]
For a fixed $P$, $ZP$ is the closest matrix to $Z$ whose rows lie in
$\operatorname{col}(P)$.

\begin{theorem}[Best approximation of prescribed rank]\label{thm:12:lowrank}
Let $Z\in\R^{n\times p}$ have positive singular values
$\sigma_1\geqslant\cdots\geqslant\sigma_r>0$ and SVD
$Z=\sum_{i=1}^r\sigma_i u_iv_i^\top$.
For $0\leqslant k\leqslant r$, put
$Z_k=\sum_{i=1}^k\sigma_i u_iv_i^\top$ (the empty sum is zero). Then
\[
\min_{\operatorname{rank}B\leqslant k}\norm{Z-B}_F^2
=\norm{Z-Z_k}_F^2=\sum_{i=k+1}^r\sigma_i^2.
\]
For $k\geqslant r$ the minimum is zero, attained by $B=Z$.
\end{theorem}
\begin{proof}
Let $B$ have row space of dimension $l\leqslant k$ and let $P$ project
onto that space. Since $B=BP$,
\[
Z-B=Z(I-P)+(ZP-B).
\]
The rows of the two summands belong to perpendicular subspaces, so
\[
\norm{Z-B}_F^2=\norm{Z(I-P)}_F^2+\norm{ZP-B}_F^2
\geqslant\norm Z_F^2-\tr(PZ^\top Z).
\]
The trace theorem applied to $Z^\top Z$ gives
$\tr(PZ^\top Z)\leqslant\sum_{i=1}^l\sigma_i^2
\leqslant\sum_{i=1}^k\sigma_i^2$, with zero eigenvalues appended as needed.
Thus every admissible $B$ has error at least the displayed tail sum.
Taking $P=\sum_{i=1}^k v_iv_i^\top$ gives
$ZP=Z_k$, of rank at most $k$.
The matrices $u_iv_i^\top$ are orthonormal in the Frobenius product,
because their pairwise inner products are
$(u_i^\top u_j)(v_i^\top v_j)=\delta_{ij}$.
Therefore $\norm{Z-Z_k}_F^2=\sum_{i>k}\sigma_i^2$, attaining the lower bound.
The cases $k=0$ and $k\geqslant r$ follow directly as well.
\end{proof}

For centered data, $\sigma_i^2=(n-1)\lambda_i(S)$.
The theorem therefore proves both the global minimum among \emph{all}
rank-$k$ matrices and the PCA residual formula
\[
\norm{Z-Z_k}_F^2=(n-1)\sum_{i>k}\lambda_i(S).
\]
It does not restrict competitors in advance to matrices of the form
$ZQQ^\top$; the proof shows why an optimal competitor can be chosen in
that form. The centered low-rank fit to the original observations is
$D_k=\mathbf1\bar d^\top+Z_k$.

For an arbitrary affine subspace $a+W$ of dimension at most $k$, with
orthogonal projector $P$ onto $W$, its total squared distance from the
observations is
\[
\sum_{i=1}^n\norm{(I-P)(d_i-a)}^2
=\norm{Z(I-P)}_F^2+n\norm{(I-P)(\bar d-a)}^2.
\]
The cross term vanishes because the centered rows sum to zero.
For each direction space $W$, the minimum is achieved exactly when
$(I-P)a=(I-P)\bar d$, meaning the affine subspace contains the mean.
Optimizing over $W$ then gives the PCA subspace through the mean.

\section*{5. SVD, sample-space computation, and scaling}
With $Z=U_r\Sigma_rV_r^\top$,
\[
S=V_r\frac{\Sigma_r^2}{n-1}V_r^\top,\qquad
G=\frac1{n-1}ZZ^\top=U_r\frac{\Sigma_r^2}{n-1}U_r^\top.
\]
The covariance matrix $S$ is $p\times p$ and concerns variable directions;
the Gram matrix $G$ is $n\times n$ and concerns observation coordinates.
They share their positive eigenvalues, with potentially different counts
of zero eigenvalues. For a unit $u$ with $Gu=\lambda u$, $\lambda>0$,
\[
v=\frac{Z^\top u}{\sqrt{(n-1)\lambda}},\qquad
Sv=\lambda v,\qquad \norm v=1,
\qquad Zv=\sqrt{(n-1)\lambda}\,u.
\]
The norm identity follows from
$\norm{Z^\top u}^2=u^\top ZZ^\top u=(n-1)\lambda$.
This provides loadings and scores from the smaller matrix when $p$ is
large compared with $n$. The formula divides by $\sqrt\lambda$ and
therefore applies only to positive eigenvalues; null directions are
obtained separately from $\ker Z$.

For the running example, $\sigma_1=\sqrt{18}=3\sqrt2$ and
$\sigma_2=\sqrt2$; the left singular columns are
\[
u_1=\tfrac12(1,1,-1,-1)^\top,\qquad
u_2=\tfrac12(1,-1,1,-1)^\top.
\]
Thus $G$ has eigenvalues $6,2/3,0,0$ while $S$ has $6,2/3$.

Suppose all variable variances are positive and put
$D_s=\operatorname{diag}(\sqrt{s_{11}},\ldots,\sqrt{s_{pp}})$.
Standardizing each column gives $Z_s=ZD_s^{-1}$ and correlation matrix
\[
R=\frac1{n-1}Z_s^\top Z_s=D_s^{-1}SD_s^{-1}.
\]
This is a congruence, not an orthogonal similarity. It changes the
optimization problem: correlation PCA measures variance after each
variable is rescaled to variance one. If a column is constant, its
variance is zero and division by that standard deviation is undefined.
Such a column contributes no centered variation and is removed from
the standardized calculation.

\begin{example}[Scaling can change the selected subspace]
If $S=\operatorname{diag}(9,1)$, covariance PCA uniquely selects the
first coordinate as its leading one-dimensional subspace and explains
$9/10$ of total variance. Standardization gives $R=I_2$.
Every unit direction is then optimal and explains $1/2$ of the standardized
total variance. Equalized units have changed both the optimum and its
uniqueness. For the running two-variable data, the two variances are
equal, so standardization multiplies all of $S$ by the same scalar;
the directions and $9/10$ proportion remain unchanged.
\end{example}

\section*{6. Min--max and compression: a variational extension}
These formulas extend the variance arguments to interior eigenvalues
and to matrices obtained by selecting or projecting variables.
For a symmetric $A\in\R^{p\times p}$ and $x\ne0$, write
$\rho_A(x)=x^\top Ax/(x^\top x)$.

\begin{theorem}[Min--max formulas]\label{thm:12:minmax}
For eigenvalues $\lambda_1\geqslant\cdots\geqslant\lambda_p$,
\[
\lambda_k=
\max_{\dim W=k}\ \min_{0\ne x\in W}\rho_A(x)
=\min_{\dim L=p-k+1}\ \max_{0\ne x\in L}\rho_A(x).
\]
\end{theorem}
\begin{proof}
For any $k$-dimensional $W$, the dimension formula implies
$W\cap\operatorname{span}(v_k,\ldots,v_p)\ne\{0\}$, since the two
dimensions sum to $p+1$. On this intersection a nonzero vector has
Rayleigh quotient at most $\lambda_k$, so the inner minimum on $W$
is at most $\lambda_k$. For
$W=\operatorname{span}(v_1,\ldots,v_k)$, every nonzero Rayleigh quotient
is at least $\lambda_k$ and equality occurs at $v_k$.
This proves the first formula, including attainment.
For any $(p-k+1)$-dimensional $L$, its intersection with
$\operatorname{span}(v_1,\ldots,v_k)$ is nonzero. Thus its inner maximum
is at least $\lambda_k$. Choosing
$L=\operatorname{span}(v_k,\ldots,v_p)$ attains this lower bound.
\end{proof}

\begin{theorem}[Separation under compression]\label{thm:12:separation}
If $Q\in\R^{p\times r}$ has orthonormal columns and
$\mu_1\geqslant\cdots\geqslant\mu_r$ are the eigenvalues of $Q^\top AQ$,
then
\[
\lambda_{p-r+j}(A)\leqslant\mu_j\leqslant\lambda_j(A)
\qquad(1\leqslant j\leqslant r).
\]
\end{theorem}
\begin{proof}
The map $y\mapsto Qy$ preserves lengths and dimensions, and
$\rho_{Q^\top AQ}(y)=\rho_A(Qy)$.
The max--min formula for $\mu_j$ ranges over $j$-dimensional subspaces
contained in $\operatorname{col}(Q)$, a subset of the competitors for
$\lambda_j(A)$. Thus $\mu_j\leqslant\lambda_j(A)$.
The min--max formula for $\mu_j$ ranges over $(r-j+1)$-dimensional
subspaces contained in that column space. Restricting the competitors
can only increase the minimum, so it is at least
$\lambda_{p-(r-j+1)+1}(A)=\lambda_{p-r+j}(A)$.
\end{proof}

Taking columns of $Q$ to be coordinate vectors yields the interlacing
bounds for any principal submatrix. Taking $A$ to be a covariance matrix
shows how the eigenvalues of a selected-variable covariance are bounded
by those of the full covariance. Individual variances need not be
eigenvalues, even though their sum always equals the eigenvalue sum.
